\section{Đánh Giá Thực Nghiệm và Kết Quả}
\label{sec:evaluation}

Phần này đánh giá framework DeSFAM về độ chính xác thực thi ở cấp độ syscall,
chi phí hiệu năng, tỷ lệ giảm thiểu, và khả năng mở rộng trong môi trường container hóa.
Các thực nghiệm được tiến hành với các khối lượng công việc thực tế và các kịch bản tấn công
để đánh giá hiệu quả và hiệu năng thời gian chạy của DeSFAM. Các chỉ số được so sánh với các
framework dẫn đầu trong điều kiện được kiểm soát nhằm đảm bảo tính công bằng và khả năng tái
lập. Kết quả làm nổi bật tác động bảo mật và tính khả thi vận hành của DeSFAM trong các triển
khai thực tế.

\subsection{Thiết Lập Thực Nghiệm}
\label{subsec:experimental_setup}

\subsubsection{Hạ Tầng}
\label{subsubsec:infrastructure}

Các thực nghiệm được tiến hành trên một máy chủ bare-metal hiệu năng cao chạy Ubuntu~24.04~LTS,
được trang bị CPU 16 nhân, 64~GB RAM, GPU NVIDIA A100 80~GB, và SSD NVMe 512~GB. Docker
26.0 với containerd~1.7 được sử dụng để quản lý vòng đời container. Môi trường eBPF được
cấu hình với libbpf~1.3 và bpftrace~0.21. Quá trình huấn luyện mô hình được thực hiện bằng
PyTorch~2.3 với tăng tốc CUDA.

\subsubsection{Khối Lượng Công Việc Container}
\label{subsubsec:container_workloads}

Chúng tôi đánh giá DeSFAM trên sáu container đa dạng: Nginx (máy chủ web), MySQL~8.0 (cơ sở
dữ liệu quan hệ), Redis~7.2 (kho lưu trữ khóa-giá trị), Node.js~20 (runtime JavaScript),
Python~3.12 (xử lý script), và MongoDB~7.0 (cơ sở dữ liệu tài liệu)---bao gồm phục vụ web,
lưu trữ dữ liệu, và thực thi script với các hồ sơ syscall khác nhau đáng kể.

\subsubsection{Sinh Khối Lượng Công Việc}
\label{subsubsec:workload_generation}

Lưu lượng thực tế được sinh ra bằng các công cụ tiêu chuẩn công nghiệp: \texttt{wrk} và ApacheBench
(\texttt{ab}) cho tải HTTP trên Nginx; \texttt{redis-benchmark} cho Redis; và các script Python
tùy chỉnh để mô phỏng các thao tác CRUD cho MySQL và MongoDB. Mỗi khối lượng công việc được chạy
ở ba mức tải: nhàn rỗi, tải trung bình (50\% công suất), và tải đỉnh (90\% công suất), trong
các khoảng thời gian 10 phút.

\subsubsection{Mô Phỏng Tấn Công}
\label{subsubsec:attack_simulation}

Ba loại kịch bản tấn công CVE thực tế được sử dụng:
\begin{itemize}
  \item \textbf{CVE-2021-4034 (Polkit pkexec)~\cite{cve_2021_4034}:} Khai thác leo thang đặc
        quyền cục bộ thông qua một chuỗi \texttt{execve} được tạo thủ công.
  \item \textbf{CVE-2022-0847 (Dirty Pipe)~\cite{cve_2022_0847}:} Khai thác ghi đè bộ nhớ kernel
        thông qua \texttt{splice} và \texttt{pipe}.
  \item \textbf{CVE-2017-5123 (waitid exploit):} Leo thang đặc quyền kernel thông qua \texttt{waitid}
        với con trỏ không hợp lệ.
\end{itemize}

Tổng cộng 14 kịch bản tấn công CVE riêng biệt đã được thực thi trên tất cả các khối lượng công
việc container.

\subsubsection{So Sánh Với Các Hệ Thống Cơ Sở}
\label{subsubsec:baselines}

DeSFAM được đánh giá chuẩn so với năm hệ thống cơ sở: Confine~\cite{ref31_confine},
Prof-gen~\cite{ref36_profgen}, NIMOS~\cite{ref38_nimos}, Optimus~\cite{ref27_optimus}, và
KubeRosy~\cite{ref29_kuberosy}. Mỗi hệ thống cơ sở được triển khai lại theo cấu hình gốc của
nó và tích hợp với Docker/Kubernetes để đảm bảo tương đồng trong khả năng thực thi thời gian
chạy và lập hồ sơ syscall.

\subsection{Các Chỉ Số Đánh Giá}
\label{subsec:metrics}

\subsubsection{Tỷ Lệ Giảm Thiểu Syscall (SRR)}
SRR định lượng khả năng của DeSFAM trong việc giảm thiểu phơi nhiễm kernel bằng cách tối thiểu
hóa các syscall được phép:
\begin{equation}
  \text{SRR} = \left(1 - \frac{|\text{DeSFAM Access List}|}{|\text{Docker Default Access List}|}\right) \times 100\%
  \label{eq:srr}
\end{equation}

Mặc định của Docker cho phép 414 syscall. Việc lập hồ sơ lai của DeSFAM (tĩnh, động,
có thông tin CVE) tạo ra các danh sách truy cập riêng cho từng container, đạt SRR cao hơn so với
các phương pháp chỉ dùng tĩnh hoặc chỉ dùng động.

\subsubsection{Hiệu Năng Phát Hiện Bất Thường}
Chúng tôi đánh giá độ chính xác của SyscallAD bằng:
\begin{itemize}
  \item Recall (TPR) -- Tỷ lệ các cuộc tấn công được nhận diện đúng
  \item Precision -- Tỷ lệ các chuỗi được gắn cờ thực sự là mối đe dọa
  \item F1-Score, AUC, và Average Precision (AP) -- Độ bền trong dữ liệu mất cân bằng
\end{itemize}
Tập hợp VAE + iForest của DeSFAM vượt trội hơn phát hiện chỉ dùng LSTM (ví dụ: NIMOS), với
độ chính xác được cải thiện và giảm dương tính giả.

\subsubsection{Độ Trễ Phát Hiện Tấn Công}
Thời gian (tính bằng $\mu$s) giữa lúc thực thi syscall độc hại và thực thi ngăn chặn. Các hook
eBPF-LSM của DeSFAM hỗ trợ chặn nội tuyến với độ trễ dưới mili giây. So với Optimus và
KubeRosy, vốn dựa trên các thành phần không gian người dùng, DeSFAM phản ứng nhanh hơn nhờ
thiết kế ở cấp độ kernel.

\subsubsection{Chi Phí Thực Thi}
Chúng tôi đánh giá mức sử dụng CPU trung bình (\%), mức tiêu thụ bộ nhớ (MB), và độ trễ xử lý
syscall ($\mu$s).

\subsubsection{Các Chỉ Số Khả Năng Thích Ứng và So Sánh}
Để đánh giá khả năng thích ứng chính sách động, chúng tôi đo lường:
\begin{itemize}
  \item Độ trễ cập nhật quy tắc qua SyscallAD hoặc các kích hoạt CVE (tính bằng $\mu$s)
  \item Tác động đến thời gian hoạt động (không quan sát thấy thời gian ngừng hoạt động)
\end{itemize}
Không có hệ thống cơ sở nào---Confine, Prof-gen, NIMOS, Optimus, hoặc KubeRosy---hỗ trợ
tự tăng cứng chính sách thời gian thực. Để so sánh toàn diện, chúng tôi đánh giá: kích thước
danh sách truy cập và SRR, tỷ lệ thành công giảm thiểu tấn công, thực thi không lỗi cho các
khối lượng công việc lành tính, và FPR trung bình cùng độ trễ thực thi.

\subsection{Kết Quả So Sánh}
\label{subsec:comparative_results}

Chúng tôi so sánh DeSFAM với năm framework lọc syscall hiện đại nhất---Confine,
Prof-gen, NIMOS, Optimus, và KubeRosy---sử dụng các chỉ số như tỷ lệ giảm thiểu syscall (SRR),
tính ổn định thực thi, ngăn chặn tấn công, và hiệu năng thời gian chạy. Đánh giá bao gồm sáu
khối lượng công việc container đa dạng: Nginx, MySQL, Redis, Node.js, Python, và MongoDB.

\subsubsection{So Sánh Tỷ Lệ Giảm Thiểu Syscall (SRR)}

Mục tiêu chính của DeSFAM là tối thiểu hóa bề mặt tấn công kernel bằng cách giảm các syscall
được phơi nhiễm. Hồ sơ seccomp mặc định của Docker cho phép 414 syscall, nhiều trong số đó là
không cần thiết và có thể bị khai thác. DeSFAM sử dụng giảm thiểu lai---phân tích tĩnh, lập
hồ sơ động, và lọc có thông tin CVE---để tạo ra các danh sách truy cập riêng cho từng container.
SRR được tính như sau:
\begin{equation}
  \text{SRR} = \left(1 - \frac{\text{Access List Size}}{414}\right) \times 100\%
  \label{eq:srr_comparative}
\end{equation}

Hình~\ref{fig:srr_bar} và Hình~\ref{fig:srr_rate} minh họa sự cải thiện trên tất cả các container
thử nghiệm. Bảng~\ref{tab:srr} tóm tắt kết quả theo từng framework.

\begin{figure}[htbp]
  \centering
  \includegraphics[width=\linewidth]{figure9}
  \caption{So Sánh Kích Thước Danh Sách Truy Cập Syscall và Tỷ Lệ Giảm Thiểu.}
  \label{fig:srr_bar}
\end{figure}

\begin{figure}[htbp]
  \centering
  \includegraphics[width=\linewidth]{figure10}
  \caption{So sánh tỷ lệ giảm thiểu syscall giữa các framework.}
  \label{fig:srr_rate}
\end{figure}

\begin{table*}[tp]
  \centering
  \caption{So Sánh Tỷ Lệ Giảm Thiểu Syscall (SRR).}
  \label{tab:srr}
  \includegraphics[width=\linewidth]{table4}
\end{table*}

DeSFAM đạt SRR trung bình 78,60\%---cao nhất trong tất cả các framework---trong khi vẫn bảo
toàn đầy đủ chức năng container. Kích thước danh sách truy cập trung bình là 89 syscall, nhỏ
nhất trong tất cả các hệ thống được so sánh.

\subsubsection{Tính Ổn Định Chức Năng}
Cả sáu container mục tiêu duy trì tính khả dụng 100\% trong suốt tất cả các khối lượng công
việc thử nghiệm. Không quan sát thấy hồi quy chức năng nào, xác nhận rằng tổng hợp danh sách
truy cập của DeSFAM cân bằng bảo mật và khả năng tương thích tốt hơn so với các bộ tạo danh
sách truy cập chỉ dùng tĩnh hoặc dựa trên nhật ký.

\subsubsection{Ngăn Chặn Tấn Công}

Khả năng chặn thời gian chạy của DeSFAM được đánh giá dựa trên một tập hợp các CVE thực tế
đại diện cho các lỗ hổng kernel Linux đa dạng. Các CVE này bao gồm nhiều loại tấn công, bao
gồm leo thang đặc quyền, thoát container, lạm dụng namespace, điều kiện đua, và khai thác
làm hỏng bộ nhớ. Để đảm bảo an toàn và khả năng tái lập, các hành vi khai thác được mô phỏng
bằng cách phát lại các chuỗi gọi syscall được ghi lại liên quan đến từng CVE, dựa trên các
thông tin công khai và phân tích bằng chứng khái niệm.

Bảng~\ref{tab:cve_blocking} tóm tắt tỷ lệ thành công chặn trên tất cả các framework. DeSFAM
đã chặn thành công tất cả các kịch bản tấn công trong thời gian thực bằng cách sử dụng phát
hiện bất thường chuỗi SyscallAD và thực thi cấp độ kernel eBPF-LSM. Ngược lại, các framework
cơ sở hoặc không phát hiện được các tấn công này hoặc thiếu khả năng thực thi đầy đủ.

\begin{table*}[tp]
  \centering
  \caption{Tỷ Lệ Thành Công Chặn Tấn Công CVE.}
  \label{tab:cve_blocking}
  \includegraphics[width=\linewidth]{table5}
\end{table*}

\subsubsection{Các Chỉ Số So Sánh Hệ Thống Cơ Sở}

\begin{table*}[tp]
  \centering
  \caption{So Sánh Các Chỉ Số Bảo Mật và Hiệu Năng với Hệ Thống Cơ Sở.}
  \label{tab:baseline_comparison}
  \includegraphics[width=\linewidth]{table6}
\end{table*}

DeSFAM đạt kích thước danh sách truy cập nhỏ nhất (89~syscall), ngăn chặn CVE đầy đủ (94\%
dấu vết CVE bị chặn), tính ổn định lành tính cao (99\%), FPR thấp nhất (1,6\%), và độ trễ
thực thi nhanh nhất (0,540~$\mu$s).

\subsection{Hiệu Năng Thời Gian Chạy}
\label{subsec:runtime_performance}

\subsubsection{Chi Phí CPU và Bộ Nhớ}

Chúng tôi đánh giá mức sử dụng tài nguyên trên sáu container đại diện có và không có DeSFAM.
Như được hiển thị trong Hình~\ref{fig:cpu_overhead} và Hình~\ref{fig:mem_overhead}, DeSFAM chỉ
tăng thêm 1,74\% mức sử dụng CPU trung bình và giảm mức tiêu thụ bộ nhớ 6,1\%, chủ yếu do
các đường dẫn syscall bị chặn và giảm thiểu tình trạng phình to bộ nhớ do khai thác gây ra.

\begin{figure}[htbp]
  \centering
  \includegraphics[width=\linewidth]{figure11}
  \caption{Mức sử dụng CPU có và không có DeSFAM. DeSFAM chỉ tạo ra chi phí CPU tối thiểu
           ($\approx$1,74\%) trong khi thực thi các chính sách syscall thời gian thực.}
  \label{fig:cpu_overhead}
\end{figure}

\begin{figure}[htbp]
  \centering
  \includegraphics[width=\linewidth]{figure12}
  \caption{Mức sử dụng bộ nhớ có và không có DeSFAM. Tiết kiệm bộ nhớ nhẹ ($\approx$6,1\%)
           là kết quả từ việc ngăn chặn syscall và giảm thiểu khai thác.}
  \label{fig:mem_overhead}
\end{figure}

\begin{table*}[tp]
  \centering
  \caption{Mức Sử Dụng CPU và Bộ Nhớ Có và Không Có DeSFAM.}
  \label{tab:cpu_memory}
  \includegraphics[width=\linewidth]{table7}
\end{table*}

Bảng~\ref{tab:overhead_summary} tóm tắt chi phí hiệu năng tổng hợp, xác nhận dấu chân hiệu
năng nhẹ của DeSFAM với mức tiết kiệm bộ nhớ trung bình 50--55~MB.

\begin{table*}[tp]
  \centering
  \caption{Tóm Tắt Chi Phí Hiệu Năng.}
  \label{tab:overhead_summary}
  \includegraphics[width=\linewidth]{table8}
\end{table*}

\subsubsection{Độ Trễ Cập Nhật Chính Sách}
Để xác nhận khả năng phản hồi của thực thi thích ứng của DeSFAM, chúng tôi đo lường thời gian
chèn và lan truyền quy tắc trong Giai đoạn~3. Khi phát hiện bất thường, các quy tắc mới được
chèn vào eBPF map và áp dụng theo thời gian thực. Độ trễ được đo bằng dấu thời gian kernel
độ phân giải cao tại cả hai giai đoạn chèn và lan truyền.

Hình~\ref{fig:policy_latency} trình bày độ trễ chèn và lan truyền quy tắc được đo trên sáu
container đang hoạt động. Bảng~\ref{tab:policy_latency} liệt kê các giá trị theo từng container.
Các giá trị được thu thập bằng dấu thời gian kernel độ phân giải cao trước và sau khi cập nhật
BPF map và lan truyền hook LSM. Tất cả các giá trị vẫn nhất quán ở mức thấp (65--82~$\mu$s),
xác nhận khả năng thời gian thực của engine thực thi trong kernel của DeSFAM. Tất cả các cập
nhật hoàn thành trong vòng 100~$\mu$s mà không làm gián đoạn dịch vụ, chứng minh khả năng
thích ứng chính sách syscall thời gian thực, không gián đoạn hoạt động của DeSFAM.

\begin{figure}[htbp]
  \centering
  \includegraphics[width=\linewidth]{figure13}
  \caption{Độ trễ cập nhật chính sách Giai đoạn~3 của DeSFAM. Tất cả quy tắc được chèn và
           thực thi trong vòng 100~$\mu$s mà không làm gián đoạn container.}
  \label{fig:policy_latency}
\end{figure}

\begin{table*}[tp]
  \centering
  \caption{Độ Trễ Cập Nhật và Lan Truyền Chính Sách.}
  \label{tab:policy_latency}
  \includegraphics[width=\linewidth]{table9}
\end{table*}

\subsubsection{Chi Phí Độ Trễ Syscall}

Chúng tôi lập hồ sơ độ trễ syscall bằng cách sử dụng histogram \texttt{bpftrace} cho các thao
tác thường xuyên như \texttt{read}, \texttt{write}, và \texttt{epoll\_wait}. Trên tất cả các
container, DeSFAM duy trì độ trễ syscall dưới mili giây, xác nhận tính phù hợp của nó cho các
môi trường thông lượng cao.

\begin{table}[htbp]
  \centering
  \caption{Chi Phí Độ Trễ Syscall (ns) cho Các Thao Tác Thông Thường.}
  \label{tab:syscall_latency}
  \begin{tabular}{lccc}
    \toprule
    \textbf{Syscall} & \textbf{Không có DeSFAM} & \textbf{Có DeSFAM} & \textbf{Delta} \\
    \midrule
    \texttt{read}       & 142\,ns & 218\,ns & +76\,ns \\
    \texttt{write}      & 156\,ns & 234\,ns & +78\,ns \\
    \texttt{epoll\_wait} & 89\,ns  & 148\,ns & +59\,ns \\
    \texttt{accept4}    & 201\,ns & 287\,ns & +86\,ns \\
    \bottomrule
  \end{tabular}
\end{table}

Bảng~\ref{tab:syscall_counts} tóm tắt các syscall có tần suất cao nhất trên mỗi container
hoàn thành trong khoảng 512--1000~$\mu$s. Hình~\ref{fig:syscall_latency_dist} trực quan hóa
số lượng tổng hợp theo từng container trong phạm vi dưới mili giây này. Các kết quả này chứng
minh khả năng của DeSFAM trong việc duy trì nhất quán độ trễ syscall thấp, ngay cả đối với các
thao tác thường xuyên như \texttt{read}, \texttt{write}, và \texttt{epoll\_wait} mà không gây
ra độ trễ đo được.

\begin{table*}[tp]
  \centering
  \caption{Số Lượng Syscall trong Phạm Vi Độ Trễ 512--1000\,$\mu$s.}
  \label{tab:syscall_counts}
  \includegraphics[width=\linewidth]{table10}
\end{table*}

\begin{figure}[htbp]
  \centering
  \includegraphics[width=\linewidth]{figure14}
  \caption{Phân phối độ trễ syscall trong phạm vi 512--1000~$\mu$s. Tất cả các đường dẫn syscall
           vẫn hoạt động hiệu quả dưới DeSFAM.}
  \label{fig:syscall_latency_dist}
\end{figure}

DeSFAM duy trì độ trễ syscall dưới mili giây cho tất cả các thao tác được đo, với delta tối
đa 86~ns cho \texttt{accept4}. Tất cả các giá trị vẫn thấp hơn nhiều so với mục tiêu thiết
kế 300~ns.

\subsubsection{Hiệu Năng Phát Hiện Bất Thường}

Mô hình SyscallAD của DeSFAM đạt được các chỉ số phân loại mạnh. Bộ phân loại tập hợp đạt
độ chính xác tổng thể 96\%, với F1-score là 0,92 trên các chuỗi tấn công. Hình~\ref{fig:roc_pr}
hiển thị các đường cong ROC và PR, với AUC là 0,94 và AP là 0,87.

\begin{figure}[htbp]
  \centering
  \includegraphics[width=\linewidth]{figure15}
  \caption{Hiệu năng mô hình SyscallAD. Trái: Đường cong Precision-Recall (AP~= 0,87). Phải:
           Đường cong ROC (AUC~= 0,94).}
  \label{fig:roc_pr}
\end{figure}

\begin{table*}[tp]
  \centering
  \caption{Các Chỉ Số Phân Loại SyscallAD.}
  \label{tab:detection_performance}
  \includegraphics[width=\linewidth]{table11}
\end{table*}

SyscallAD đạt độ chính xác tổng thể 96\%, F1-score là 0,92, và AUC là 0,94 trên tập dữ liệu
DongTing. Điều này thể hiện sự cải thiện đáng kể so với các phương pháp chỉ dùng VAE (F1~= 0,84)
hoặc chỉ dùng iForest (F1~= 0,79), xác nhận lợi ích của chiến lược hợp nhất tập hợp.

\subsubsection{So Sánh Tính Năng với Các Phương Pháp Hiện Có}
DeSFAM tích hợp độc đáo phát hiện bất thường thời gian thực, thực thi chính sách thích ứng,
và tối thiểu hóa syscall trong một framework thống nhất. Bảng~\ref{tab:feature_comparison} nêu
bật ưu điểm của nó so với các framework hiện có trên 12~chiều kỹ thuật.

\begin{table*}[tp]
  \centering
  \caption{So Sánh Tính Năng với Các Framework Hiện Có.}
  \label{tab:feature_comparison}
  \includegraphics[width=\linewidth]{table12}
\end{table*}

\subsection{Nghiên Cứu Loại Bỏ: Đóng Góp Theo Từng Giai Đoạn}
\label{subsec:ablation}

Để hiểu tác động của từng giai đoạn trong DeSFAM, chúng tôi tiến hành nghiên cứu loại bỏ vô
hiệu hóa từng thành phần một. Như được hiển thị trong Bảng~\ref{tab:ablation}, Giai đoạn~1
đơn lẻ đạt tỷ lệ giảm thiểu syscall 62,4\% nhưng chỉ chặn được 3 trong số 14 CVE. Việc thêm
Giai đoạn~2 (SyscallAD) tăng cường phát hiện bất thường, nâng số lượng bị chặn lên 8 và giảm
đáng kể dương tính giả. Giai đoạn~3 cung cấp thực thi cấp độ kernel, chặn 10 CVE ngay cả khi
không có hỗ trợ ML. Khi cả ba giai đoạn đều hoạt động, DeSFAM chặn tất cả 14 CVE và đạt tỷ
lệ giảm thiểu syscall 81,6\%.

\begin{table}[htbp]
  \centering
  \caption{Nghiên Cứu Loại Bỏ: Hiệu Ứng Theo Từng Giai Đoạn đối với Bảo Mật.}
  \label{tab:ablation}
  \includegraphics[width=\linewidth]{table13}
\end{table}

Nghiên cứu này xác nhận rằng tất cả các giai đoạn đều thiết yếu: Giai đoạn~1 tối thiểu hóa
phơi nhiễm, Giai đoạn~2 tăng cường phát hiện, và Giai đoạn~3 đảm bảo thực thi bền vững.
